\paragraph{Neutral completions and linear source types.}
Call $X$ a \emph{source sink} when its passive status affects no other
complete source predicate. For a source-reconstructing operation $T$ define
\[
 \mathsf S_XT=T+\partial_XT,
\]
regarding the residual as an operation ignoring $X$. A zero residual
contributes no summand. This is a smaller completion than adjoining the
union of all other input menus.

\begin{proposition}[Least sink completions]\label{core:prop:sink-completions}
The map $\mathsf S_X$ is a monotone closure operator on core operations,
and $\mathsf S_XT$ is the least core majorant of $T$ making $X$ a source
sink. It preserves $X$'s complete source predicate and replaces every
other predicate by its restriction to neutral status at $X$.
The closures commute. For $A\subseteq V$, their simultaneous completion is
\[
 \mathsf S_AT=\bigcup_{B\subseteq A}\partial_BT.
\]
Completing all sources gives input union on the essential support.
For nonempty $V$ one also has the pointwise identity
$T=\bigcap_{X\in V}\mathsf S_XT$.
\end{proposition}
\begin{proof}
The ghost comparison gives $T\subseteq X\cup\partial_XT$.
At any other source's probe, neutralizing $X$ can only enlarge its
coefficient, so the union with the residual gives precisely that neutral
restriction. At an $X$-source probe the residual contributes no queried
winner, leaving that predicate unchanged. Thus the completion is a sink.
Conversely a sink operation $G$ contains its neutral residual: its older
source contributions are unchanged by $X$. If $G\supseteq T$, then
$G\supseteq\partial_XG\supseteq\partial_XT$, proving minimality.
Extensivity and monotonicity follow from the formula, and minimality
gives idempotence.

Neutral residuals commute and distribute over union, winner by winner.
Expanding the commuting closures gives the displayed formula.
The one-input residual at an essential source is its projection, while
an inessential source contributes zero. Provenance gives the reverse
containment in their input union.

Finally suppose $k\notin T(v)$. If some input $X$ contains $k$,
replacing it by $I\setminus\{k\}$ changes no foreign-$k$ atomic flag
and removes only that input's $k$-contribution. Source decomposition
therefore gives $k\notin\partial_XT(v)$. If no input contains $k$,
no residual introduces it either. In both cases some sink completion
still omits $k$, proving the intersection identity.
\end{proof}

\begin{lemma}[Linearizing a marked occurrence]
\label{core:lem:neutral-marked-linearization}
For every marked occurrence $o$ of $X$ in a core word $T(X,v)$ there is
a primitive spine $S_o(X,v)$ with exactly one $X$-occurrence such that
\[
 S_o\subseteq T+\partial_XT.
\]
Its full owner-factor type equals that of $o$: for every owner, the
conjunction of negative sibling profiles, evaluated with all $X$-flags
false, is unchanged.
\end{lemma}
\begin{proof}
For a subword $U$ write $C_U=\partial_XU$, possibly zero, and let its
deleted $r$-profile be its $r$-profile with $X$ non-$r$.
The residual compiler supplies three facts:
\begin{enumerate}
\item $U\subseteq X\cup C_U$;
\item every nonzero deleted profile is preserved exactly in $C_U$;
\item if $U=\{r\}$ but $C_U\ne\{r\}$, then $X=\{r\}$.
\end{enumerate}
The last follows from singleton-profile factorization and the second
fact when $X$ is not the $r$ singleton; a zero deleted profile cannot
produce that singleton in this case.

Follow the marked path from its leaf. At an $r$-owned node, orient the
marked child as $T_1$ and its sibling as $T_2$. If $T_2$'s deleted
$r$-profile is zero, skip the node. Otherwise attach $C_2$ by
$\star_r$. This is an $X$-free core helper, and its profile is exact.
Thus every original marked owner factor is retained: skipped nodes
had trivial factors, and attached helpers have the same nonzero profiles.
The word stays a primitive spine with one $X$.

For containment, assume its child spine satisfies $S\subseteq T_1+C_1$.
Let $a,b$ indicate whether the two deleted $r$-profiles are nonzero,
and let $C$ be the parent residual from the four-row compiler.
If $b=0$, the discard rules and the compiler give $C_1\subseteq C$.
An alleged winner of the skipped spine outside the parent value and $C$
would belong to $T_1\setminus C_1\subseteq X$ and be killed by an
$r$-absorbing sibling $T_2$. It is foreign to $r$, whereas $b=0$ and
$T_2=\{r\}$ force $X=\{r\}$, a contradiction.

Suppose $b=1$, so the new spine is $S\star_r C_2$.
If $C_2=\{r\}$, the new spine and the parent residual are both $\{r\}$.
If $C_2\ne\{r\}$ but $a=1$ and $C_1=\{r\}$, profile preservation
gives $T_1=S=\{r\}$, so the new spine is again safe.
In every remaining case the compiler gives $C=C_1\cup C_2$.
Any unsafe winner must therefore come from $T_1\setminus C_1\subseteq X$.
The root cannot lose its own owner, so such a winner is foreign to $r$;
losing it requires $T_2=\{r\}$. Since $C_2\ne\{r\}$, fact~(3)
again forces $X=\{r\}$, a contradiction.
Empty residues obey the discard rules; $b=1$ ensures the attached
$C_2$ is nonempty, and $C_1=0$ fits the last union case.
This completes the safety induction.
\end{proof}

\begin{corollary}[Linear covers of arbitrary source intervals]
\label{core:cor:interval-linear-cover}
Every core $F\in\mathcal I(C)$ is the union of $C$ and primitive
one-$X$ spines bounded by $F$. Every actual full $X$-occurrence type
of any realizer of $F$ is realized exactly by one such spine.
More generally these conclusions hold for $\mathsf S_XT$ with respect
to every marked type in $T$, for arbitrary core $T$.
\end{corollary}
\begin{proof}
The interval bounds give $\partial_XF=C$, so the preceding spines
lie below $F+C=F$. A winner of $F$ outside $C$ must come from an
$X$-occurrence. Its preserved owner factors make the corresponding
spine supply it too. Union with $C$ therefore recovers $F$.
For arbitrary $T$, union $\partial_XT$ with the spines obtained from
its marked occurrences. They lie below $T+\partial_XT$ and cover
every original $X$-source contribution. The remaining contributions
are contained in the residual, since neutralizing $X$ only enlarges
the other source coefficients. This gives the final assertion.
\end{proof}

The neutral completion is essential in the marked lemma: an arbitrary
core operation can require repeated occurrences even for a single positive
source probe. The interval bound makes exactly those neutral contributions
available to protect the linear spine.

\begin{lemma}[Profile inheritance along a linear spine]
\label{core:lem:linear-profile-inheritance}
For $K\geq3$, let $C$ be any core baseline, and let a primitive one-$X$ spine $S$ with
passive helpers satisfy $S\subseteq C\cup X$. Write $B_r$ for its full
owner factors, $\psi_r=\neg B_r$, and $R=\{r:\psi_r\ne0\}$.
Then:
\begin{enumerate}
\item $\Phi_r(C)\leq\psi_r$ for every active owner $r\in R$;
\item every helper $H$ satisfies $\Phi_a(C)\leq\Phi_a(H)$ for every
  inactive owner $a\notin R$;
\item every helper satisfies $H(w)\subseteq C(w)$ whenever
  $\psi_r(E_r(w))=0$ for all owners $r$.
\end{enumerate}
\end{lemma}
\begin{proof}
For each owner,
$\psi_r=\bigvee_{\text{$r$-owned path gates}}\Phi_r(H)$.
If $r$ is active, the deleted $r$-profile of the spine is nonzero:
on the all-$r$ passive diagonal its first $r$-gate absorbs, and all later
helpers remain $r$. The residual compiler preserves this profile.
Since $\partial_XS\subseteq C$ and the deleted profile recurrence is
bounded by the displayed disjunction,
\[
 \Phi_r(C)\leq\Phi_r(\partial_XS)\leq\psi_r.
\]

An inactive owner $a$ has no path gate, since every nonempty helper
has a positive top profile. When $C=\{a\}$ and $X=\{a\}$,
the bound forces $S=\{a\}$. Its $a$-singleton predicate is the
conjunction of all helper $a$-profiles, giving~(2).

For~(3), every helper avoids its gate owner's singleton at $w$.
Give $X$ the ambiguous menu $I\setminus\{k\}$, for any $k\notin C(w)$.
The prefix stays ambiguous, every gate takes union, and the bound makes
every helper omit $k$. Varying $k$ gives the asserted containment.
\end{proof}

\begin{corollary}[Inactive-owner boxes]
\label{core:cor:inactive-owner-box}
Let $F\in\mathcal I(C)$ retain its distinguished source on the
all-$a$ passive diagonal. Then
\[
 \theta_a(E):=\bigwedge_{r\ne a}\Phi_a(C)(W\setminus E_r)
 \quad\Longrightarrow\quad P_F(E)=1
\]
for every proper coloring $E$.
\end{corollary}
\begin{proof}
Choose an $X$-occurrence supplying a foreign winner at that diagonal and
linearize it. Its factor $B_a(W)=1$ makes $a$ inactive. Every helper
therefore inherits $\Phi_a(C)$. If an $r$-helper were the $r$ singleton
at $E_r$ with $r\ne a$, its positive $r$-coalition would be disjoint
from the inherited positive $a$-coalition $W\setminus E_r$, violating
profile admissibility. Thus every owner factor of the spine is true at
$E$. Its sound source box is contained in $P_F$, proving the implication.
\end{proof}

\begin{corollary}[Canonical sink completions need not glue]
\label{core:cor:sink-completion-obstruction}
For every fixed $K\geq3$ there are noncore operations $F$ at unbounded
arity with complete source-dependence graph, such that every nonempty
simultaneous sink completion $\mathsf S_AF$ is core.
Each single-source completion $\mathsf S_XF$ has acyclic dependence;
if $|A|\geq2$, $\mathsf S_AF$ is ordinary input union.
These canonical completions commute and satisfy
$F=\bigcap_X\mathsf S_XF$ pointwise.
\end{corollary}
\begin{proof}
Use the prime Gray operation from \cref{core:cor:prime-obstructions}.
Its proper neutral residual at $X$ is input union on the other inputs.
Thus $\mathsf S_XF$ is exactly the greatest completion of its single
source family, retaining every other input, and is core by
\cref{core:thm:source-completion}. Those other source predicates are
unconditional, so its dependence is acyclic. Two distinct neutral
residuals already have full input union as their union; this proves the
assertion for larger $A$. The closure formulas and intersection proof
apply to the source-reconstructing Gray operation as well.

Every input $i$ is the free coordinate of some Gray region, since each
bit flips in both directions along the cyclic walk. Put $\{a\}$ at
that coordinate, with $a\notin\{b,c\}$, and the prescribed $b/c$
atoms elsewhere. This is a canonical source-$i$ probe whose output omits $a$.
Changing any other input $j$ to $I\setminus\{a\}$ creates two free
coordinates, so no exceptional region remains and $a$ is present.
Hence $\pi_i$ depends on every $j\ne i$, proving completeness of the
source graph. Noncoreness is the Gray obstruction theorem.
\end{proof}
